\cotefile{33}
\scanpage{50}
\mspage{\unsure{21}.}
\uline{Le \unsure{thm} de dualité pour une immersion}.

\margin{La \unsure{dém.} \unsure{s'applique} \unsure{pour} un \unsure{morphisme} \uline{étale} \unsure{également}}

Soit $i:Y\longrightarrow X$ une immersion.
Alors $i_!$ est exact (en fait, identique à $i_*$ si
$i$ est une imm. fermée), donc
$\mathbb{R}i_!$ est de nature \unsure{triviale}
\[
\mathbb{R}i_!=i_!:\quad \mathrm{D}^+(Y)\longrightarrow\mathrm{D}^+(X)
\]
(\unsure{inutile} ici de prendre des résolutions \unsure{injectives})

D'autre part, on définit un foncteur
\[
i^!:\deleted{\unsure{i}}\quad \mathrm{C}(Y)\longrightarrow\mathrm{C}(X)
\]
de la façon suivante : si $i$ est un
composé
\diagram{\begin{tikzcd}[ampersand replacement=\&]
Y \arrow[r,no head,"i'"] \& U \arrow[r,no head,"j"] \& X
\end{tikzcd}}
, avec $i'$ une
imm. \uline{fermée}, $j$ une imm. \unsure{ouverte},
et on prend :
\[
i^!(F)=\deleted{\unsure{\mathbb{R}}}i'^*\underline{\Gamma}_Y(j^*(F))
\]
Ce foncteur ne dépend pas, à isom \unsure{unique}
près, du choix de la factorisation choisie,
comme on voit \unsure{p.ex.} grâce à la
formule d'adjonction
\[
\operatorname{Hom}_Y(G,i^!(F))\simeq\operatorname{Hom}_X(i_!(G),F)
\]
qui dit que $i^!$ \deleted{est adjoint à droite de $i_!$}.
Comme \unsure{il|la} \deleted{\unsure{foncteur|formule}} \unsure{existe|évidente} un \unsure{isom.}
canonique de cette \unsure{nature} pour \deleted{\unsure{i}}
l'imm. \unsure{ouverte} $j$ \unsure{entre} $j_!$ et $j^*$ (\unsure{fait}
\mspage{49}
